\chapter{Conclusions and Future Work}\label{chp:conclusion}

\needspace{5\baselineskip}
\section{Summary of the Work}

This research was designed to find out if ML-based runtime predictions could increase job scheduling effectiveness in the use of large-scale GPU clusters. A major issue with current cluster scheduling is there exists a basic informational gap between the time a user submits their jobs to be processed and when those jobs are actually completed. Thus, a scheduler has no way of approximating SJF, which has been proven to result in minimum average job wait times among non-preemptive policies. This informational gap results in increased queueing delay, lower utilization of available hardware resources, and waste of energy.

Therefore, this thesis developed and tested a complete end-to-end pipeline based on the Alibaba PAI GPU cluster traces~\cite{weng2022mlaas}. The pipeline consisted of three primary modules:
\begin{enumerate}
    \item Data cleaning and feature generation (including a new sweep-line-based cluster utilization metric);
    \item Training and comparative analysis of six predictive models across two frameworks; and
    \item An event-driven scheduling simulator to translate differences in prediction accuracy into differences in scheduling outcomes.
\end{enumerate}

\needspace{5\baselineskip}
\section{Key Findings}

The four major results we obtained with these experiments are as follows:

\textbf{Scheduling Duality (Production vs.\ Prediction):} Although XGBoost and tree ensemble methods produced the smallest point-prediction errors (MAE) and were ranked first overall, our DL model, specifically the Categorical LSTM, performed best for true scheduling performance at \emph{both} cluster scales tested (32 GPUs and 256 GPUs), not only at the larger one.\footnote{Exact JCT-reduction percentages are being refreshed following a code-level correction to the hyperparameter-tuning loss function and the addition of non-learned reference baselines; see the Limitations section.} Therefore this demonstrates that achieving an accurate ordinal rank order (with or without taking into account the presence of many micro-jobs) has much greater impact in reducing the HoL blocking problem than does having high-quality regression predictions -- a conclusion that, on the present evidence, does not depend on cluster size.

\textbf{The Sequence Trap:} All sliding-window temporal sequence models performed worse than their static counterparts. Additionally, removing categorical features from sliding-window models resulted in catastrophic failures. Specifically, SJF-LSTM (Numeric Sequence) exhibited a degradation in scheduling performance relative to a blind FIFO baseline of $-$45.41\%. Furthermore, injecting random temporal variance into the input data while providing no stable embedding representations for the categorical entities resulted in unstable queues.

\textbf{Feature Informativeness of Sweep-Line Utilization:} Our concurrent job count feature at submission time, determined via sweep-line algorithm, is one of the top predictors in MDI analysis across all tree-based models. As such this supports our prior decision to include dynamic cluster load information in the input feature space.

\textbf{ML-SJF substantially improves upon baselines:} For each of the static ML models used, ML-SJF outperformed both FIFO and SRF heuristics by significant margins in terms of mean job wait time and P95 (tail) wait times. Thus we have evidence for how great the marginal benefit of obtaining exact job duration predictions is beyond just knowing job resource requirements.

\needspace{5\baselineskip}
\section{Contributions}

These are some of the contributions to the field of cluster resource management made through this thesis:

\begin{enumerate}

    \item \textbf{End-to-end prediction-to-scheduling pipeline using production trace.} All elements from data collection, feature generation, model training, and scheduling simulation have been integrated into a single reproducible system. By integrating these elements together, we can exactly map prediction performance to scheduling results; something that cannot be accomplished using just one component.

    \item \textbf{Sweep-line cluster utilization feature.} A new sweep-line cluster utilization feature based upon the current utilization of the cluster at submission time of each job was created. It is consistently found to be among the top predictive features and frequently omitted from existing GPU cluster prediction research.

    \item \textbf{Cross-paradigm controlled comparative analysis of different models.} Six prediction models were trained and evaluated across two paradigms within the same experimentally controlled environment and dataset. No cross-paradigm controlled comparative analysis exists in the literature today and provides unambiguous guidance for practitioners seeking to choose models for similar problems.

    \item \textbf{Conversion of raw point-prediction accuracy to actual scheduling outcomes.} An empirical analysis of the relationship between raw point-prediction MAE and actual scheduling efficiency showed that a model's ordinal ranking quality, not its point-prediction accuracy, is what determines scheduling outcomes, and that this relationship holds consistently at both cluster scales evaluated rather than being specific to large-scale deployments.

    \item \textbf{Reproducible open artifacts.} The preprocessing, model training and simulation code have been documented and released to allow other researchers to modify or reproduce the experiments presented in this thesis.

    \clearpage
    \item \textbf{Industrial applicability.} The ML-SJF framework is intended to replace user-supplied unreliable wall-time estimates for any SJF-capable scheduler. Therefore it is directly applicable to any GPU cluster manager (e.g., Kubernetes with GPU device plugins), HPC scheduler (e.g., SLURM's Backfill Algorithm) or cloud batch queueing system (e.g., YARN's Capacity Scheduler).

\end{enumerate}

\needspace{5\baselineskip}
\section{Limitations}

There were several areas of limitation in regards to this research:

\textbf{Offline Prediction Setting.} All prediction models in this thesis were developed using historical data, and they are run when the user submits their request. Therefore, these models are not updated as new data is submitted. If we had a live production system with changing workload distributions (concept drift) then we would have needed to retrain our models periodically to keep the predictions accurate.

\textbf{No Preemption.} The simulator did not allow for job preemption: once a job started running, it ran until it was completed. However, many production systems do use preemption for high-priority jobs. Adding such capabilities to the simulator would greatly increase its complexity.

\textbf{Single-Node Job Placement.} The simulator placed every job on one node. Although distributed training jobs can take place across multiple nodes, most of the jobs in the Alibaba traces were single-node, so we feel comfortable that this simplification does not limit our ability to generalize the results to multi-node workloads.

\textbf{Limitations from Feature Availability.} To get the ``concurrent'' job count features, I needed to know how long all previous jobs took to complete. Unfortunately, when making decisions in a live environment we will not know this information in advance. Instead, we will need to make an estimate of how long a given job will take based upon what has been estimated previously, which adds another layer of estimation and thus another potential source of error.

\textbf{Results Under Revision.} A post-defense scientific review of this codebase identified a mismatch between the loss function used during hyperparameter search and the one used to fit the final LightGBM/XGBoost models, and found that no non-learned reference baseline (a per-user or per-job-profile median lookup) was reported alongside the tuned models. Both are being corrected: the final-refit loss function now matches the search objective, and two non-learned baselines (Per-User Median, Profile Median) have been added to the prediction and scheduling tables. Several specific percentages reported in this chapter (JCT improvements, Oracle-capture ratios) were computed before these corrections and are being refreshed accordingly; the qualitative finding that ranking quality, not point-prediction accuracy, drives scheduling gains -- and that this holds at both cluster scales tested -- was re-verified against the current checkpoints and is unaffected.

\needspace{5\baselineskip}
\section{Threats to Validity}

\textbf{External Validity.} All experiments were run using a single production trace from the Alibaba PAI cluster. Although the trace is one of the largest publicly available traces from a GPU cluster workload, the results may not generalize to other types of clusters due to differences in hardware components (for example, CPU-only HPC clusters), users or workload distribution. The relative ranking of tree-based and DL models found here is specific to the Alibaba trace's heavy-tailed, predominantly ML workload characteristics and may not transfer to workloads with a different runtime distribution.

\textbf{Internal Validity.} No outlier filtering was applied to job runtimes prior to simulation; the heavy-tailed runtime distribution described in Chapter~\ref{chp:dataset} (including a small number of extremely long-running jobs) is therefore fully reflected in the reported mean-based statistics. Because the identical workload trace and job ordering were replayed for every scheduling policy, this does not bias the relative comparison between policies, but it does mean that a small number of extreme jobs can have an outsized effect on mean-based aggregates such as mean JCT and mean waiting time. The reported median and tail-percentile statistics are less sensitive to this effect and should be consulted alongside the means when interpreting policy differences.

\textbf{Construct Validity.} The simulator treats actual job execution times as being determinate and knowable at completion time. In practice, job execution times can be affected by several dynamic factors including network contention, memory constraints, etc. None of these factors are modeled. Furthermore, the lack of job preemption, job migration and multi-node gang scheduling limits the complexity of the resource model compared to what actually occurs in real-world production systems. This is an acceptable tradeoff when simulating scheduler behavior given common practices in the literature, and while limiting the direct applicability of simulated statistics to actual deployment statistics, they still serve as useful estimates.

\textbf{Statistical Conclusion Validity.} Each pairwise comparison of scheduling policies is substantiated through a Wilcoxon signed-rank test. Nevertheless, this thesis did not calculate bootstrap confidence intervals for the mean JCT improvement associated with each policy pair. This would have provided additional insight into the uncertainty surrounding the reported point estimate. Authors wishing to enhance statistical rigor should incorporate confidence interval estimation via bootstrapping into their future research efforts.

\needspace{5\baselineskip}
\section{Future Work}

The findings of this thesis present many avenues for future study:

\textbf{Decision-Focused Learning from Start to Finish.} Although this thesis has shown that there is no direct link between the accuracy of a model's predictions of job completion times and its performance in terms of scheduling quality, future research should attempt to train a full decision-making system, including both the predictive component and the sorting function, using a single end-to-end trainable architecture. The inclusion of differentiable sorting or reinforcement learning will enable the creation of a complete decision-focused learning pipeline from simulation output.

\textbf{Job Scheduling via Learning to Rank.} As the scheduler can rely solely upon knowledge of the relative order of jobs in the queue to perform its task, it would seem reasonable to cast the problem of predicting runtime as an L2R problem in the future. Using pairwise or listwise ranking loss functions (such as RankNet or ListMLE) will provide a means to formally maximize the capacity of the scheduler to protect ``mice'' from ``elephants''. Moreover, future evaluations of predictive schedulers will need to utilize ranking correlation measures (e.g., Kendall's $\tau$, Spearman's $\rho$, or NDCG) as primary metrics, eliminating the disconnect between statistical regression metrics and the final sorting goal, thereby providing a theoretically valid basis for predictive scheduling.

\textbf{Per-User Sequential Modeling.} Future studies can address the ``sequence trap'' that is present within globally shared cluster queues by creating user-specific and session-specific sequences. By grouping a user's history in chronological order before providing these groups as sequential inputs to an embedding algorithm, recurrent neural networks (e.g., LSTM) can be used to model user-specific submission patterns. This approach avoids the temporal disruptions caused by the interleaved submissions of other users.

\textbf{Online and Adaptive Prediction.} Developing the pipeline into an online environment in which incremental retraining occurs as each new record becomes available would lead to increased resilience against distribution shifts. Techniques such as online gradient boosting or continual neural network learning may be employed toward achieving this goal.

\textbf{Preemptive Scheduling.} Adding the capability for the simulator to include preemption will allow for the exploration of even more complex scheduling strategies such as SRPT, which is optimally suited to minimize average waiting time under conditions where preemption is allowed.

\textbf{Support for Distributed Jobs.} Extending the resource model to allow for jobs that require resources on multiple machines and/or subject to gang-scheduling constraints (i.e., all required machines must be concurrently available) would result in a simulation tool capable of simulating distributed DL training workflows.

\textbf{Prediction with Uncertainty.} Current models predict job runtimes as point estimates. Quantile regression or conformal prediction may be used to obtain probabilistic bounds on predicted job runtimes; these bounds could then be utilized by a risk-aware scheduler to make either more conservative or more aggressive scheduling decisions based upon the priority preferences of the user.

\textbf{Cross-Trace Generalization.} Evaluation of the predictive and scheduling pipeline using trace data from clusters owned by other organizations (such as Microsoft or Google) will evaluate whether the findings generalize across differing machine architectures and workload mixes.